\documentclass[letterpaper, 12pt]{article}

% ============================================================
% PAQUETES (Estandarizados)
% ============================================================
\usepackage[utf8]{inputenc}
\usepackage[spanish]{babel}
\usepackage{amsmath, amssymb, amsfonts}
\usepackage{geometry}
\usepackage{fancyhdr}
\usepackage{enumitem}
\usepackage{graphicx}
\usepackage{hyperref}
\usepackage{multicol}
\usepackage{parskip}
\usepackage{xcolor}
\usepackage{tcolorbox}
\usepackage{eso-pic}
\usepackage{transparent}
\usepackage{tikz}
\tcbuselibrary{skins, breakable}

% ============================================================
% GEOMETRÍA Y COLORES
% ============================================================
\geometry{letterpaper, margin=1.5cm, top=3cm, bottom=2.5cm, headheight=2cm}

\definecolor{celesteSuave}{HTML}{F0F7FF}
\definecolor{azulFuerte}{HTML}{1A5276}
\definecolor{turquesaUsach}{HTML}{33CCCC}
\definecolor{enlace}{HTML}{1A5276}

% Rutas de imágenes (Asegúrate de tener la carpeta /images)
\newcommand{\logoup}{./imagenes/logo_up.png}
\newcommand{\logousach}{./imagenes/logo_usach.png}
\newcommand{\logofondo}{./imagenes/logo_up.png}

% ============================================================
% INFORMACIÓN DE LA GUÍA (Francisca Varela - Álgebra II)
% ============================================================
\newcommand{\institucion}{Usach Premium}
\newcommand{\curso}{Álgebra II}
\newcommand{\guiasnumero}{Guía: Sistemas de Ecuaciones y Matrices}
\newcommand{\semestre}{1er. Semestre de 2026}
\newcommand{\preparadopor}{Francisca Varela}
\newcommand{\webinstitucion}{www.usachpremium.cl}
\newcommand{\instagraminstitucion}{https://www.instagram.com/usach.premium}
\newcommand{\transparencia}{0.05}
\newcommand{\fraseportada}{"Por y para estudiantes"}

\newcommand{\listacontenidos}{
    \item[\color{azulFuerte}$\Rightarrow$] Determinantes y Rango
    \item[\color{azulFuerte}$\Rightarrow$] Método de Gauss-Jordan
    \item[\color{azulFuerte}$\Rightarrow$] Teorema del Rango
    \item[\color{azulFuerte}$\Rightarrow$] Ecuaciones Matriciales
}

% ============================================================
% CONFIGURACIÓN DE ENCABEZADOS Y MARCA DE AGUA
% ============================================================
\pagestyle{fancy}
\fancyhf{}
\fancyhead[L]{\includegraphics[height=1.2cm]{\logoup}}
\fancyhead[R]{%
    \begin{minipage}[b]{0.42\textwidth}
        \raggedleft\fontsize{9pt}{11pt}\selectfont
        \textbf{\color{azulFuerte}\institucion} \\
        \curso\ — \guiasnumero \\
        \textit{\color{gray}@usach.premium}
    \end{minipage}\hspace{0.1cm}%
    \includegraphics[height=1.2cm]{\logousach}%
}
\fancyfoot[L]{\fontsize{9pt}{11pt}\selectfont \textit{\fraseportada}}
\fancyfoot[R]{\fontsize{9pt}{11pt}\selectfont Página \thepage}

\newcommand{\MarcaAgua}{
    \AddToShipoutPicture{\AtPageCenter{\makebox(0,0){%
        \transparent{\transparencia}\includegraphics[width=0.7\textwidth]{\logofondo}%
    }}}
}

% Entornos de Cajas
\newenvironment{kitpropiedades}[1]{%
    \begin{tcolorbox}[colback=celesteSuave, colframe=azulFuerte, title=\textbf{#1}, fonttitle=\bfseries]
}{\end{tcolorbox}}

\newenvironment{solbox}[1]{%
    \begin{tcolorbox}[colback=celesteSuave, colframe=azulFuerte, title=\textbf{#1}, fonttitle=\bfseries, breakable]
}{\end{tcolorbox}}

% Entornos de Sección
\newenvironment{ejercicios}{\begin{center}\begin{tcolorbox}[colback=azulFuerte, colframe=azulFuerte, colupper=white, width=0.8\textwidth, center, arc=10pt]\centering\Large\textbf{EJERCICIOS}\end{tcolorbox}\end{center}\MarcaAgua\small}{}
\newenvironment{soluciones}{\newpage \begin{center}\begin{tcolorbox}[colback=azulFuerte, colframe=azulFuerte, colupper=white, width=0.8\textwidth, center, arc=10pt]\centering\Large\textbf{SOLUCIONARIO}\end{tcolorbox}\end{center}\MarcaAgua\small}{}

% ============================================================
\begin{document}

% PORTADA INSTITUCIONAL
\thispagestyle{empty}
\AddToShipoutPicture*{%
    \put(54, 700){\includegraphics[width=2cm]{\logoup}}
    \put(420, 706){\includegraphics[width=5cm]{\logousach}}
}

\vspace*{0.5cm}
\begin{center}
    \color{azulFuerte}
    {\huge \textbf{GUÍA DE EJERCICIOS}} \\[0.5cm]
    {\Huge \textbf{\curso}} \\[0.3cm]
    {\Large \guiasnumero} \\[1.5cm]

    \begin{tcolorbox}[colback=celesteSuave, colframe=azulFuerte, arc=10pt, width=0.85\textwidth, center, boxrule=1.5pt]
        \centering \vspace{0.2cm} {\Large \textbf{Contenidos}} \vspace{0.3cm}
        \color{black}
        \begin{itemize}[leftmargin=3cm] \listacontenidos \end{itemize}
        \vspace{0.2cm}
    \end{tcolorbox}
\end{center}

\vspace{4.5cm}
\begin{center}
    {\Large \textbf{\semestre}} \\[0.8cm]
    {\large \textbf{Preparado por: \preparadopor}}
\end{center}

\vspace{4cm}
\begin{center}
    {\color{azulFuerte}\hrule height 0.5pt} \vspace{0.3cm}
    \begin{minipage}{0.45\textwidth} \centering \textbf{Instagram:} \\ \small\href{\instagraminstitucion}{@usach.premium} \end{minipage}
    \begin{minipage}{0.45\textwidth} \centering \textbf{Web:} \\ \small\href{http://\webinstitucion}{\webinstitucion} \end{minipage}
\end{center}
\newpage

%%%%%%%% EJERCICIOS %%%%%%%
\begin{ejercicios}

\begin{kitpropiedades}{Recordatorio: Teorema del Rango}
    Sea $A \in M_{m \times n}(\mathbb{R})$. El sistema $Ax=b$ es:
    \begin{itemize}
        \item \textbf{Compatible Determinado:} si $rg(A) = rg(A|b) = n$ (n° de variables).
        \item \textbf{Compatible Indeterminado:} si $rg(A) = rg(A|b) < n$.
        \item \textbf{Incompatible:} si $rg(A) \neq rg(A|b)$.
    \end{itemize}
\end{kitpropiedades}


\begin{enumerate}[label=\textbf{\arabic*.}]
% --- EJERCICIO 1 ---
    \item Calcule el determinante de la matriz, escalónela, redúzcala y determine su rango:
    \[
    A = \begin{pmatrix}
    1 & 2 & -1 & 3 \\
    2 & 4 & 1 & 7 \\
    -1 & -2 & 2 & -1 \\
    3 & 6 & 0 & 8
    \end{pmatrix}
    \]
% --- EJERCICIO 2 ---
    \item Resuelva el siguiente sistema de ecuaciones utilizando el método de Gauss-Jordan:
    \[
    \begin{cases}
    x + 2y - z = 1 \\
    2x + 4y + z = 5 \\
    -x - 2y + 2z = 0
    \end{cases}
    \]
% --- EJERCICIO 3 ---
    \item Dado el sistema de ecuaciones homogéneo:
    \[
    \begin{cases}
    x + 2y - z = 0 \\
    2x + 4y - 2z = 0 \\
    3x + 6y - 3z = 0
    \end{cases}
    \]
    Utilizando el Teorema del rango, indique si tiene solución única o infinitas soluciones. En caso de ser infinitas, describa el conjunto solución.
% --- EJERCICIO 4 ---
    \item Considere las matrices simétricas $U, W, T, Z \in M_{\mathbb{R}}(n)$, tales que $|U + T| \neq 0$. 
    Demuestre que:
    \[ \text{Si } XT - W^t = (Z - UX^t)^t \quad \text{entonces} \quad X = (W + Z)(U + T)^{-1} \]
% --- EJERCICIO 5 ---
    \item Dado el sistema:
    \[
    \begin{cases}
    x + y + z = 2 \\
    2x + 3y + 4z = 5 \\
    x + 2y + (a + 1)z = 0
    \end{cases}
    \]
    Determine los siguientes conjuntos:
    \[
    \begin{aligned}
    S_1 &= \{a \in \mathbb{R} / \text{el sistema tiene solución única}\} \\
    S_2 &= \{a \in \mathbb{R} / \text{el sistema tiene infinitas soluciones}\} \\
    S_3 &= \{a \in \mathbb{R} / \text{el sistema no tiene solución}\}
    \end{aligned}
    \]
% --- EJERCICIO 6 ---
    \item Sea $B \in M_{\mathbb{R}}(4)$, determine su inversa, si es que existe:
    \[
    B = \begin{pmatrix}
    2 & 1 & -3 & -2 \\
    4 & 2 & 0 & 0 \\
    1 & -2 & 3 & 4 \\
    0 & 3 & 1 & 1
    \end{pmatrix}
    \]
% --- EJERCICIO 7 ---
    \item Dado el sistema:
    \[
    \begin{cases}
    x + y - z = 1 \\
    2x + 2y - 2z = 2 \\
    x + 2y + az = 3 \\
    3x + 4y + (a-2)z = 5
    \end{cases}
    \]
    Determine el conjunto: 
    \[ S = \{ a \in \mathbb{R} \ / \ \text{el sistema es compatible} \} \]
% --- EJERCICIO 8 ---
    \item Considerando el sistema:
    \[
    \begin{cases}
    \begin{aligned}
    x + y + z &= 1 \\
    2x + 2y + (a + b)z &= 2 \\
    x + y + az &= b
    \end{aligned}
    \end{cases}
    \]
    Determine los siguientes conjuntos:
    \[
    \begin{aligned}
    S_1 &= \{(a, b) \in \mathbb{R}^2 / \text{el sistema tiene solución única}\} \\
    S_2 &= \{(a, b) \in \mathbb{R}^2 / \text{el sistema tiene infinitas soluciones}\} \\
    S_3 &= \{(a, b) \in \mathbb{R}^2 / \text{el sistema no tiene solución}\}
    \end{aligned}
    \]
% --- EJERCICIO 9 ---
    \item Considere la matriz:
    \[
    A = \begin{pmatrix}
    1 & 2 & a \\
    0 & 1 & 3 \\
    2 & 4 & a+1
    \end{pmatrix}
    \]
    Determine los valores de $a \in \mathbb{R}$ para los cuales $A$ es invertible.
% --- EJERCICIO 10 ---
    \item Utilizando el Teorema del rango, determine $k \in \mathbb{R}$ tal que el siguiente sistema sea incompatible:
    \[
    \begin{cases}
    x + 2y + z + 2w = 1 \\
    -2x - 3y - z - (k+4)w = k-1 \\
    x + y + z + (k+1)w = 2-k \\
    2x + 3y + 2z + (5-k)w = k^2 - 2
    \end{cases}
    \]
% --- EJERCICIO 11 ---
    \item Sean $A, B \in M_{\mathbb{R}}(2)$ tales que:
    \[
    A = \begin{pmatrix} 2 & 1 \\ 1 & 1 \end{pmatrix}, B = \begin{pmatrix} 3 & 1 \\ 2 & 0 \end{pmatrix}
    \]
    Resuelva la ecuación matricial:
    \[
    AX = B
    \]
% --- EJERCICIO 12 ---
    \item Dado el sistema:
    \[
    \left\{ \begin{aligned}
    x + 2y - z + w &= 1 \\
    2x + 4y - 2z + 2w &= 2 \\
    x + 3y + az + (a - 1)w &= 2 \\
    3x + 7y + (a - 2)z + (a + 1)w &= 4
    \end{aligned} \right.
    \]
    \begin{enumerate}
        \item[a)] Utilizando el Teorema del rango, determine los valores de $a \in \mathbb{R}$ para los cuales el sistema es compatible.
        \item[b)] Resuelva el sistema para $a = 2$.
    \end{enumerate}
% --- EJERCICIO 13 ---
    \item Dado el sistema:
    \[
    \begin{cases}
    \begin{aligned}
    x + 3y + (\beta - 1)z &= 4 \\
    \alpha x + 9y + (\alpha\beta - \alpha)z &= 4\alpha \\
    -x + (6 - 3\alpha)y + (4 - \beta)z &= 2 \\
    2x + 6y + (2\beta - 5)z &= 11 - \alpha^2
    \end{aligned}
    \end{cases}
    \]
    Utilizando el Teorema del rango, determine el conjunto:
    \[
    S = \{(\alpha, \beta) \in \mathbb{R}^2 / \text{el S.E.L tiene infinitas soluciones}\}
    \]
% --- EJERCICIO 14 ---
    \item Determine el conjunto solución asociado al sistema:
    \[
    \begin{cases}
    x_1 + 2x_2 + 3x_3 = 2 \\
    3x_1 + 6x_2 + x_3 + 8x_4 = -3 \\
    5x_1 + 10x_2 - x_3 + 17x_4 = 12
    \end{cases}
    \]
\end{enumerate}
\end{ejercicios}

%%%%%%%%%% SOLUCIONARIO %%%%%%%%%
\begin{soluciones}
% --- EJERCICIO 1 ---
\begin{solbox}{Ejercicio 1}
\begin{itemize}
    \item \textbf{Calculamos el determinante utilizando el método de Laplace + Sarrus.}
\end{itemize}
\textbf{Paso 1: Laplace}
\[|A| = \begin{vmatrix} 1 & 2 & -1 & 3 \\ 2 & 4 & 1 & 7 \\ -1 & -2 & 2 & -1 \\ 3 & 6 & 0 & 8 \end{vmatrix} = 1 \begin{vmatrix} 4 & 1 & 7 \\ -2 & 2 & -1 \\ 6 & 0 & 8 \end{vmatrix} - 2 \begin{vmatrix} 2 & 1 & 7 \\ -1 & 2 & -1 \\ 3 & 0 & 8 \end{vmatrix} + (-1) \begin{vmatrix} 2 & 4 & 7 \\ -1 & -2 & -1 \\ 3 & 6 & 8 \end{vmatrix} - 3 \begin{vmatrix} 2 & 4 & 1 \\ -1 & -2 & 2 \\ 3 & 6 & 0 \end{vmatrix}\]
\textbf{Paso 2: Sarrus}
\begin{gather*}
\begin{vmatrix} 4 & 1 & 7 \\ -2 & 2 & -1 \\ 6 & 0 & 8 \end{vmatrix} = -10; \quad \begin{vmatrix} 2 & 1 & 7 \\ -1 & 2 & -1 \\ 3 & 0 & 8 \end{vmatrix} = -5; \quad \begin{vmatrix} 2 & 4 & 7 \\ -1 & -2 & -1 \\ 3 & 6 & 8 \end{vmatrix} = 0; \quad \begin{vmatrix} 2 & 4 & 1 \\ -1 & -2 & 2 \\ 3 & 6 & 0 \end{vmatrix} = 0 \\
|A| = 1 \cdot (-10) - 2 \cdot (-5) + (-1) \cdot 0 - 3 \cdot 0 = -10 + 10 + 0 + 0 = 0 \implies \boxed{|A| = 0}
\end{gather*}
\begin{center}
\fbox{
\begin{minipage}{0.9\textwidth}
    \textbf{Observación:} El determinante de una matriz es cero si existen filas o columnas linealmente dependientes. En este caso, como $Columna_2 = 2 \cdot Columna_1$, se tiene que $|A| = 0$, sin necesidad de cálculos adicionales.
\end{minipage}
}
\end{center}

 \begin{itemize}
    \item \textbf{Escalonamos la matriz utilizando operaciones de línea.}
\end{itemize}
\[
A =\left( \begin{array}{cccc} 1 & 2 & -1 & 3 \\ 2 & 4 & 1 & 7 \\ -1 & -2 & 2 & -1 \\ 3 & 6 & 0 & 8 \end{array} \right) \xrightarrow{\substack{R_2 \to R_2 - 2R_1 \\ R_3 \to R_3 + R_1 \\ R_4 \to R_4 - 3R_1}} \left( \begin{array}{cccc} 1 & 2 & -1 & 3 \\ 0 & 0 & 3 & 1 \\ 0 & 0 & 1 & 2 \\ 0 & 0 & 3 & -1 \end{array} \right) \xrightarrow{R_2 \leftrightarrow R_3} \left( \begin{array}{cccc} 1 & 2 & -1 & 3 \\ 0 & 0 & 1 & 2 \\ 0 & 0 & 3 & 1 \\ 0 & 0 & 3 & -1 \end{array} \right) \]

\[ \xrightarrow{\substack{R_3 \to R_3 - 3R_2 \\ R_4 \to R_4 - 3R_2}} \left( \begin{array}{cccc} 1 & 2 & -1 & 3 \\ 0 & 0 & 1 & 2 \\ 0 & 0 & 0 & -5 \\ 0 & 0 & 0 & -7 \end{array} \right) \xrightarrow{R_4 \to R_4 - \frac{7}{5}R_3} \left( \begin{array}{cccc} 1 & 2 & -1 & 3 \\ 0 & 0 & 1 & 2 \\ 0 & 0 & 0 & -5 \\ 0 & 0 & 0 & 0 \end{array} \right)
\]
\begin{itemize}
    \item \textbf{Reducimos para obtener su forma escalonada reducida.}
\end{itemize}
\[
\begin{pmatrix} 
1 & 2 & -1 & 3 \\ 
0 & 0 & 1 & 2 \\ 
0 & 0 & 0 & -5 \\ 
0 & 0 & 0 & 0 
\end{pmatrix} 
\xrightarrow{R_3 \to \left(\frac{-1}{5}\right) \cdot R_3} 
\begin{pmatrix} 
1 & 2 & -1 & 3 \\ 
0 & 0 & 1 & 2 \\ 
0 & 0 & 0 & 1 \\ 
0 & 0 & 0 & 0 
\end{pmatrix}
\xrightarrow{\substack{R_1 \to R_1 - 3R_3 \\ R_2 \to R_2 - 3R_3}}
\begin{pmatrix} 
1 & 2 & -1 & 0 \\ 
0 & 0 & 1 & 0 \\ 
0 & 0 & 0 & 1 \\ 
0 & 0 & 0 & 0 
\end{pmatrix}
\]
\[
\xrightarrow{R_1 \to R_1 + R_2} 
\begin{pmatrix} 
1 & 2 & 0 & 0 \\ 
0 & 0 & 1 & 0 \\ 
0 & 0 & 0 & 1 \\ 
0 & 0 & 0 & 0 
\end{pmatrix}
\]
\begin{center}
    \fbox{Tenenmos 3 pivotes o escalones, por lo tanto el $rg(A) = 3$.}
\end{center}
\end{solbox}
% --- EJERCICIO 2 ---
\begin{solbox}{Ejercicio 2}
    \begin{itemize}
    \item \textbf{Aplicamos el método Gauss-Jordan}
\end{itemize}
% Primera parte del escalonamiento
\[
\begin{cases} 
    x + 2y - z = 1 \\ 
    2x + 4y + z = 5 \\ 
    -x - 2y + 2z = 0 
\end{cases} 
\color{red}\rightarrow\color{black}
\left( \begin{array}{ccc|c} 
    1 & 2 & -1 & 1 \\ 
    2 & 4 & 1 & 5 \\ 
    -1 & -2 & 2 & 0 
\end{array} \right)
\xrightarrow{\substack{R_2 \to R_2 - 2R_1 \\ R_3 \to R_3 - 3R_1}}
\left( \begin{array}{ccc|c} 
    1 & 2 & -1 & 1 \\ 
    0 & 0 & 3 & 3 \\ 
    0 & 0 & 1 & 1 
\end{array}\right)
\]
\[
\xrightarrow{R_2 \leftrightarrow R_3}
\left( \begin{array}{ccc|c} 
    1 & 2 & -1 & 1 \\ 
    0 & 0 & 1 & 1 \\ 
    0 & 0 & 3 & 3 
\end{array} \right)
\xrightarrow{R_3 \to R_3 - 3 \cdot R_2}
\left( \begin{array}{ccc|c} 
    1 & 2 & -1 & 1 \\ 
    0 & 0 & 1 & 1 \\ 
    0 & 0 & 0 & 0 
\end{array} \right)
\xrightarrow{R_1 \to R_1 + R_2}
\left( \begin{array}{ccc|c} 
    1 & 2 & 0 & 2 \\ 
    0 & 0 & 1 & 1 \\ 
    0 & 0 & 0 & 0 
\end{array} \right)
\]
\[
\begin{cases} 
    x + 2y = 2 \\ 
    z = 1 
\end{cases}
\color{red}\rightarrow\color{black}
\begin{cases} 
    x = 2 - 2y \\ 
    z = 1 
\end{cases}
\]
\begin{center}
    \fbox{\textbf{Conjunto solución:} $(x, y, z) = (2-2t, t, 1)$ con $t \in \mathbb{R}$.}
\end{center}
\end{solbox}
% --- EJERCICIO 3 ---
\begin{solbox}{Ejercicio 3}
    \begin{itemize}
    \item \textbf{Aplicamos el método Gauss-Jordan al sistema homogéneo.}
\end{itemize}
\[
\begin{cases} 
    x + 2y - z = 0 \\ 
    2x + 4y - 2z = 0 \\ 
    3x + 6y - 3z = 0 
\end{cases} 
\color{red}\rightarrow\color{black}
\left( \begin{array}{ccc|c} 
    1 & 2 & -1 & 0 \\ 
    2 & 4 & -2 & 0 \\ 
    3 & 6 & -3 & 0 
\end{array} \right)
\xrightarrow{\substack{R_2 \to R_2 - 2R_1 \\ R_3 \to R_3 - 3R_1}} 
\left( \begin{array}{ccc|c} 
    1 & 2 & -1 & 1 \\ 
    0 & 0 & 0 & 0 \\ 
    0 & 0 & 0 & 0 
\end{array}\right)
\]

Como $rg(A) = rg(A|B) = 1 < Ca = 3$, el sistema tiene infinitas soluciones.\\
Para describir las infinitas soluciones, partimos de la única ecuación que queda: $ x + 2y -z = 0$ \\
Tomamos $y = s$ y $z = t$, entonces: $x = -2y + z \implies x = -2s + t$
\begin{center}
    \fbox{\textbf{Conjunto solución:} $(x, y, z) = (-2s + t, s, t)$ con $s, t \in \mathbb{R}$.}
\end{center}
\end{solbox}
% --- EJERCICIO 4 ---
\begin{solbox}{Ejercicio 4}
    \begin{itemize}
    \item \textbf{Utilizaremos las propiedades de la transpuesta y la hipótesis de que las matrices son simétricas, lo que implica que $M^t = M$ para $M \in \{U, W, T, Z\}$}
\end{itemize}

\textbf{Desarrollamos el lado derecho de la igualdad:}
\begin{enumerate}
    \item \textbf{Aplicamos la propiedad $(A - B)^t = A^t - B^t$ y $(AB)^t = B^t A^t$:}
    \[ (Z - UX^t)^t = Z^t - (UX^t)^t = Z^t - (X^t)^t U^t \]
    
    \item \textbf{Aplicamos la propiedad de simetría ($Z^t = Z, U^t = U$) y la doble transpuesta ($(X^t)^t = X$):}
    \[ (Z - UX^t)^t = Z - XU \]
    
    \item \textbf{Sustituimos en la ecuación original considerando que $W^t = W$:}
    \[ XT - W = Z - XU \]
    
    \item \textbf{Agrupamos los términos con $X$ a un lado y las matrices constantes al otro:}
    \[ XT + XU = Z + W \]
    
    \item \textbf{Factorizamos $X$ por la izquierda:}
    \[ X(T + U) = Z + W \]
    
    \item \textbf{Como $|U + T| \neq 0$, la matriz $(T + U)$ es invertible. Multiplicamos por $(T + U)^{-1}$ por la derecha:}
    \[ X = (W + Z)(U + T)^{-1} \]
\end{enumerate}

\begin{center}
    \fbox{Queda demostrado que $X = (W + Z)(U + T)^{-1}$}
\end{center}
\end{solbox}
% --- EJERCICIO 5 ---
\begin{solbox}{Ejercicio 5}
    \begin{itemize}
    \item \textbf{Aplicamos el método Gauss-Jordan:}
\end{itemize}
\[
\begin{cases} 
    x + y + z = 2 \\ 
    2x + 3y + 4z = 5 \\ 
    x + 2y + (a + 1)z = 0 
\end{cases} 
\color{red}\rightarrow\color{black}
\left( \begin{array}{ccc|c} 
    1 & 1 & 1 & 2 \\ 
    2 & 3 & 4 & 5 \\ 
    1 & 2 & a+1 & 0 
\end{array} \right)
\xrightarrow{\substack{R_2 \to R_2 - 2R_1 \\ R_3 \to R_3 - R_1}} 
\left( \begin{array}{ccc|c} 
    1 & 1 & 1 & 2 \\ 
    0 & 1 & 2 & 1 \\ 
    0 & 1 & a & -2 
\end{array} \right)
\]
\[
\xrightarrow{R_3 \to R_3 - R_2}
\left( \begin{array}{ccc|c} 
    1 & 1 & 1 & 2 \\ 
    0 & 1 & 2 & 1 \\ 
    0 & 0 & a-2 & -3 
\end{array} \right)
\quad \text{La matriz está en su forma \textbf{escalonada}}
\]
\begin{itemize}
    \item \textbf{Ahora veremos por casos y aplicaremos el Teorema del Rango:}
\end{itemize}

\textbf{Caso 1: } $a \neq 2$
\[
\text{Rg}(A) = 3, \quad \text{Rg}(A|B) = 3, \quad \text{Cant. Variables} = 3 \implies \text{Solución única}
\]

\textbf{Caso 2: } $a = 2$
\[
\text{Rg}(A) = 2, \quad \text{Rg}(A|B) = 3 \implies \text{El sistema \textbf{no posee solución}}
\]

\begin{center}
\fbox{
\begin{minipage}{0.55\textwidth}
    \textbf{Conclusión:} \\
    $S_1 = \{a \in \mathbb{R} / \text{ el sistema tiene solución única}\} = \mathbb{R} - \{2\}$ \\
    $S_2 = \{a \in \mathbb{R} / \text{ el sistema tiene infinitas soluciones}\} = \emptyset$ \\
    $S_3 = \{a \in \mathbb{R} / \text{ el sistema no tiene solución}\} = \{2\}$
\end{minipage}
}
\end{center}
\end{solbox}

\newpage
% --- EJERCICIO 6 ---
\begin{solbox}{Ejercicio 6}
    \begin{itemize}
    \item \textbf{Análisis previo:} Para que $B^{-1}$ exista, se debe cumplir que $|B| \neq 0$. 
    Calculando el determinante obtenemos que:
    \[ |B| = 82 \]
    \begin{center}
    Como $|B| \neq 0$, la matriz es \textbf{invertible}.
    \end{center}
\end{itemize}

\begin{itemize}
    \item \textbf{Calculamos la inversa de la matriz $B \in M_{\mathbb{R}}(4)$ mediante el método de Gauss-Jordan sobre $[B \mid I]$:}
\end{itemize}
\[
\left( \begin{array}{cccc|cccc} 
2 & 1 & -3 & -2 & 1 & 0 & 0 & 0 \\ 
4 & 2 & 0 & 0 & 0 & 1 & 0 & 0 \\ 
1 & -2 & 3 & 4 & 0 & 0 & 1 & 0 \\ 
0 & 3 & 1 & 1 & 0 & 0 & 0 & 1 
\end{array} \right)
\xrightarrow{R_1 \leftrightarrow R_3}
\left( \begin{array}{cccc|cccc} 
1 & -2 & 3 & 4 & 0 & 0 & 1 & 0 \\ 
4 & 2 & 0 & 0 & 0 & 1 & 0 & 0 \\ 
2 & 1 & -3 & -2 & 1 & 0 & 0 & 0 \\ 
0 & 3 & 1 & 1 & 0 & 0 & 0 & 1 
\end{array} \right)
\]

\[
\xrightarrow{\substack{R_2 \to R_2 - 4R_1 \\ R_3 \to R_3 - 2R_1}} 
\left( \begin{array}{cccc|cccc} 
1 & -2 & 3 & 4 & 0 & 0 & 1 & 0 \\ 
0 & 10 & -12 & -16 & 0 & 1 & -4 & 0 \\ 
0 & 5 & -9 & -10 & 1 & 0 & -2 & 0 \\ 
0 & 3 & 1 & 1 & 0 & 0 & 0 & 1 
\end{array} \right)
\xrightarrow{R_2 \to \frac{1}{10}R_2}
\left( \begin{array}{cccc|cccc} 
1 & -2 & 3 & 4 & 0 & 0 & 1 & 0 \\ 
0 & 1 & -1.2 & -1.6 & 0 & 0.1 & -0.4 & 0 \\ 
0 & 5 & -9 & -10 & 1 & 0 & -2 & 0 \\ 
0 & 3 & 1 & 1 & 0 & 0 & 0 & 1 
\end{array} \right)
\]

\begin{center}
    \textit{(Aplicando operaciones elementales de fila para reducir a la identidad. Son muchas, por lo tanto el orden y tener paciencia son claves.)}
\end{center}
\[
\xrightarrow{\text{Reducción completa}}
\left( \begin{array}{cccc|cccc} 
    1 & 0 & 0 & 0 & -1/41 & 21/82 & 1/41 & -6/41 \\ 
    0 & 1 & 0 & 0 & 2/41 & -1/82 & -2/41 & 12/41 \\ 
    0 & 0 & 1 & 0 & -29/41 & 35/82 & -12/41 & -10/41 \\ 
    0 & 0 & 0 & 1 & 23/41 & -16/41 & 18/41 & 15/41 
\end{array} \right)
\]

\begin{center}
\fbox{
\begin{minipage}{0.40\textwidth}
    \centering
    \textbf{Resultado final:} \\
    \medskip
    $B^{-1} = \frac{1}{82} \begin{pmatrix} 
        -2 & 21 & 2 & -12 \\ 
        4 & -1 & -4 & 24 \\ 
        -58 & 35 & -24 & -20 \\ 
        46 & -32 & 36 & 30 
    \end{pmatrix}$
\end{minipage}
}
\end{center}
\end{solbox}

\newpage
% --- EJERCICIO 7 ---
\begin{solbox}{Ejercicio 7}
    \begin{itemize}
    \item \textbf{Aplicamos el método Gauss-Jordan:}
\end{itemize}
\[
\begin{cases} 
    x + y - z = 1 \\ 
    2x + 2y - 2z = 2 \\ 
    x + 2y + az = 3 \\
    3x + 4y + (a-2)z = 5
\end{cases} 
\color{red}\rightarrow\color{black}
\left( \begin{array}{ccc|c} 
    1 & 1 & -1 & 1 \\ 
    2 & 2 & -2 & 2 \\ 
    1 & 2 & a & 3 \\ 
    3 & 4 & a-2 & 5 
\end{array} \right)
\xrightarrow{\substack{R_2 \to R_2 - 2R_1 \\ R_3 \to R_3 - R_1 \\ R_4 \to R_4 - 3R_1}} 
\left( \begin{array}{ccc|c} 
    1 & 1 & -1 & 1 \\ 
    0 & 0 & 0 & 0 \\ 
    0 & 1 & a+1 & 2 \\ 
    0 & 1 & a+1 & 2 
\end{array} \right)
\]
\[
\xrightarrow{R_2 \leftrightarrow R_3}
\left( \begin{array}{ccc|c} 
    1 & 1 & -1 & 1 \\ 
    0 & 1 & a+1 & 2 \\ 
    0 & 0 & 0 & 0 \\ 
    0 & 1 & a+1 & 2 
\end{array} \right)
\xrightarrow{R_4 \to R_4 - R_2}
\left( \begin{array}{ccc|c} 
    1 & 1 & -1 & 1 \\ 
    0 & 1 & a+1 & 2 \\ 
    0 & 0 & 0 & 0 \\ 
    0 & 0 & 0 & 0 
\end{array} \right)
\]
\begin{itemize}
    \item Observamos que, tras el escalonamiento, las filas 3 y 4 se anulan por completo independientemente del valor de $a$, teniendo $\text{Rg}(A) = \text{Rg}(A \mid B) = 2$.
    \item Como $\text{Rg}(A) = \text{Rg}(A \mid B)$, el sistema tiene solución (compatible) para cualquier valor de $a \in \mathbb{R}$.
\end{itemize}

\begin{center}
\fbox{
\begin{minipage}{0.6\textwidth}
    \centering
    \textbf{Resultado:} \\
    \medskip
    $S = \{a \in \mathbb{R} \mid \text{el sistema es compatible}\} = \mathbb{R}$
\end{minipage}
}
\end{center}
\end{solbox}
% --- EJERCICIO 8 ---
\begin{solbox}{Ejercicio 8}
    \begin{itemize}
    \item \textbf{Escalonamos la matriz aumentada $(A \mid B)$ para identificar los pivotes:}
\end{itemize}
\[
(A \mid B) = \left( \begin{array}{ccc|c} 
    1 & 1 & 1 & 1 \\ 
    2 & 2 & a+b & 2 \\ 
    1 & 1 & a & b 
\end{array} \right)
\begin{aligned}
    \xrightarrow{\substack{R_2 \to R_2 - 2R_1 \\ R_3 \to R_3 - R_1}} 
\end{aligned}
\left( \begin{array}{ccc|c} 
    1 & 1 & 1 & 1 \\ 
    0 & 0 & a+b-2 & 0 \\ 
    0 & 0 & a-1 & b-1 
\end{array} \right)
\]
\begin{itemize}
    \item \textbf{Estudio de Casos:}
\end{itemize}
Observamos que la segunda columna no tiene pivote, por lo que el sistema nunca tendrá solución única ($S_1 = \emptyset$). Analizamos los casos entre las filas 2 y 3:

\begin{enumerate}
    \item \textbf{Caso para infinitas soluciones ($S_2$):} 

Para que el sistema sea compatible indeterminado, necesitamos que el rango sea menor al número de incógnitas ($\text{Rg}(A) = \text{Rg}(A|B) < 3$). Esto ocurre si logramos que las dos últimas filas de la matriz se anulen completamente ($0 = 0$).

Revisando nuestra matriz escalonada:
\[ \left( \begin{array}{ccc|c} 
    1 & 1 & 1 & 1 \\ 
    0 & 0 & a+b-2 & 0 \\ 
    0 & 0 & a-1 & b-1 
\end{array} \right) \]

\newpage
Para que las filas 2 y 3 sean filas de ceros, deben cumplirse \textbf{simultáneamente} estas condiciones:
\begin{itemize}
    \item De la Fila 3: $a-1 = 0 \implies \mathbf{a = 1}$ y $b-1 = 0 \implies \mathbf{b = 1}$.
    \item Verificamos en Fila 2: Si $a=1$ y $b=1$, entonces $1+1-2 = 0$. ¡Se cumple!
\end{itemize}
\begin{center}
\fbox{
\begin{minipage}{0.90\textwidth}
    \textit{Nota adicional:} También habría infinitas soluciones si la Fila 2 y la Fila 3 fueran dependientes, pero dado que el lado derecho de la Fila 2 es $0$, la única forma de no generar una contradicción es que ambas filas de coeficientes se anulen o sean consistentes con $b$.
\end{minipage}
}
\end{center}

\textbf{2. Caso para ninguna solución ($S_3$):}

El sistema es incompatible si al realizar el escalonamiento aparece una contradicción lógica del tipo $0 = \text{constante} \neq 0$. Esto sucede específicamente en la Fila 3:
\begin{itemize}
    \item \textbf{Condición}: Si fijamos $\mathbf{a = 1}$, la parte izquierda de la Fila 3 se anula:
    \[ (1 - 1)z = 0z = 0 \]
    
    \item \textbf{Contradicción}: Si al mismo tiempo elegimos cualquier valor de $\mathbf{b \neq 1}$, la parte derecha de la Fila 3 será:
    \[ b - 1 \neq 0 \]
    
    \item \textbf{Resultado}: Obtenemos la ecuación $0 = b - 1$. Como $b - 1$ no es cero, esto es un \textbf{absurdo matemático}.
\end{itemize}

Por lo tanto, el sistema es incompatible siempre que $a = 1$ y $b \neq 1$.
\end{enumerate}

\begin{center}
\fbox{
\begin{minipage}{0.90\textwidth}
    \textbf{Conclusión:} \\
    $S_1 = \{(a, b) \in \mathbb{R}^2 \mid \text{el sistema tiene solución única}\} = \emptyset$ \\
    $S_2 = \{(a, b) \in \mathbb{R}^2 \mid \text{el sistema tiene infinitas soluciones}\} = \{(a, b) \in \mathbb{R}^2 \mid a=1, b=1 \text{ o } a+b=2, a \neq 1\}$ \\
    $S_3 = \{(a, b) \in \mathbb{R}^2 \mid \text{el sistema no tiene solución}\} = \{(a, b) \in \mathbb{R}^2 \mid a=1, b \neq 1\}$
\end{minipage}
}
\end{center}
\end{solbox}
% --- EJERCICIO 9 ---
\begin{solbox}{Ejercicio 9}
    \begin{enumerate}
    \item \textbf{Análisis de Invertibilidad (Determinante):}
    
Una matriz cuadrada es invertible si y solo si su determinante es diferente de cero. Calculamos $|A|$ utilizando la regla de Sarrus o expansión por la primera columna (ya que tiene un cero):

\[ |A| = \begin{vmatrix} 1 & 2 & a \\ 0 & 1 & 3 \\ 2 & 4 & a + 1 \end{vmatrix} \]

Expandiendo por la primera columna:
\[ |A| = 1 \cdot \begin{vmatrix} 1 & 3 \\ 4 & a + 1 \end{vmatrix} - 0 \cdot \begin{vmatrix} 2 & a \\ 4 & a + 1 \end{vmatrix} + 2 \cdot \begin{vmatrix} 2 & a \\ 1 & 3 \end{vmatrix} \]

\newpage
Calculamos los determinantes $2 \times 2$:
\begin{itemize}
    \item $1 \cdot [(1)(a + 1) - (3)(4)] = a + 1 - 12 = a - 11$
    \item $2 \cdot [(2)(3) - (a)(1)] = 2(6 - a) = 12 - 2a$
\end{itemize}

Sumamos los resultados:
\[ |A| = (a - 11) + (12 - 2a) \]
\[ |A| = -a + 1 \]

    \item \textbf{Condición de Invertibilidad:}

Para que $A$ sea invertible, requerimos que $|A| \neq 0$:
\[ -a + 1 \neq 0 \implies \mathbf{a \neq 1} \]
\end{enumerate}

\begin{center}
\fbox{
\begin{minipage}{0.6\textwidth}
    \centering
    \textbf{Conclusión:} \\
    \medskip
    La matriz $A$ es invertible para todo el conjunto: \\
    $S = \{a \in \mathbb{R} \mid a \neq 1\} = \mathbb{R} - \{1\}$
\end{minipage}
}
\end{center}
\end{solbox}
% --- EJERCICIO 10 ---
\begin{solbox}{Ejercicio 10}
    \begin{enumerate}
    \item \textbf{Análisis mediante el Teorema del Rango:}

Escalonamos la matriz aumentada $(A \mid B)$ mediante operaciones elementales:

\[
\left( \begin{array}{cccc|c} 
    1 & 2 & 1 & 2 & 1 \\ 
    -2 & -3 & -1 & -(k+4) & k-1 \\ 
    1 & 1 & 1 & k+1 & 2-k \\ 
    2 & 3 & 2 & 5-k & k^2-2 
\end{array} \right)
\xrightarrow{\substack{R_2 \to R_2 + 2R_1 \\ R_3 \to R_3 - R_1 \\ R_4 \to R_4 - 2R_1}}
\left( \begin{array}{cccc|c} 
    1 & 2 & 1 & 2 & 1 \\ 
    0 & 1 & 1 & -k & k+1 \\ 
    0 & -1 & 0 & k-1 & 1-k \\ 
    0 & -1 & 0 & 1-k & k^2-4 
\end{array} \right)
\]

\[
\xrightarrow{\substack{R_3 \to R_3 + R_2 \\ R_4 \to R_4 + R_2}}
\left( \begin{array}{cccc|c} 
    1 & 2 & 1 & 2 & 1 \\ 
    0 & 1 & 1 & -k & k+1 \\ 
    0 & 0 & 1 & -1 & 2 \\ 
    0 & 0 & 1 & 1-2k & k^2+k-3 
\end{array} \right)
\xrightarrow{R_4 \to R_4 - R_3}
\left( \begin{array}{cccc|c} 
    1 & 2 & 1 & 2 & 1 \\ 
    0 & 1 & 1 & -k & k+1 \\ 
    0 & 0 & 1 & -1 & 2 \\ 
    0 & 0 & 0 & 2-2k & k^2+k-5 
\end{array} \right)
\]

\item \textbf{Condición de Incompatibilidad:}

Para que el sistema sea \textbf{incompatible}, la última fila debe representar una contradicción ($0 = \text{constante} \neq 0$). Esto sucede si:
\begin{enumerate}
    \item $2 - 2k = 0 \implies \mathbf{k = 1}$
    \item $k^2 + k - 5 \neq 0$
\end{enumerate}

Sustituimos $k = 1$ en la expresión del término independiente:
\[ (1)^2 + (1) - 5 = 1 + 1 - 5 = -3 \]
Como $\mathbf{-3 \neq 0}$, se genera la contradicción $0 = -3$ cuando $k = 1$.
\end{enumerate}
\vspace{0.4cm}

\begin{center}
\fbox{
\begin{minipage}{0.5\textwidth}
    \centering
    \textbf{Resultado:} \\
    \medskip
    El sistema es incompatible para $\mathbf{k = 1}$.
\end{minipage}
}
\end{center}
\end{solbox}
% --- EJERCICIO 11 ---
\begin{solbox}{Ejercicio 11}
    \begin{enumerate}
    \item \textbf{Análisis de existencia de la solución:}

Para despejar $X$, la matriz $A$ debe ser invertible ($|A| \neq 0$). Calculamos su determinante:
\[ |A| = \begin{vmatrix} 2 & 1 \\ 1 & 1 \end{vmatrix} = (2 \cdot 1) - (1 \cdot 1) = 2 - 1 = 1 \]
Como $\mathbf{|A| = 1 \neq 0}$, la matriz $A$ es invertible y la solución es única: $\mathbf{X = A^{-1}B}$.

\item \textbf{Cálculo de $A^{-1}$:}

Para una matriz $2 \times 2$, la inversa se obtiene mediante: $A^{-1} = \frac{1}{|A|} \begin{pmatrix} d & -b \\ -c & a \end{pmatrix}$
\[ A^{-1} = \frac{1}{1} \begin{pmatrix} 1 & -1 \\ -1 & 2 \end{pmatrix} = \begin{pmatrix} 1 & -1 \\ -1 & 2 \end{pmatrix} \]

\item \textbf{Resolución de la ecuación $X = A^{-1}B$:}

\[ X = \begin{pmatrix} 1 & -1 \\ -1 & 2 \end{pmatrix} \begin{pmatrix} 3 & 1 \\ 2 & 0 \end{pmatrix} \]
Multiplicamos fila por columna:
\begin{itemize}
    \item $x_{11} = (1 \cdot 3) + (-1 \cdot 2) = 3 - 2 = 1$
    \item $x_{12} = (1 \cdot 1) + (-1 \cdot 0) = 1 + 0 = 1$
    \item $x_{21} = (-1 \cdot 3) + (2 \cdot 2) = -3 + 4 = 1$
    \item $x_{22} = (-1 \cdot 1) + (2 \cdot 0) = -1 + 0 = -1$
\end{itemize}
\end{enumerate}

\begin{center}
\fbox{
\begin{minipage}{0.2\textwidth}
    \centering
    \textbf{Resultado:} \\
    \medskip
    $X = \begin{pmatrix} 1 & 1 \\ 1 & -1 \end{pmatrix}$
\end{minipage}
}
\end{center}
\end{solbox}
% --- EJERCICIO 12 ---
\begin{solbox}{Ejercicio 12}
    \textbf{a) Análisis de compatibilidad mediante el Teorema del Rango:}

Escalonamos la matriz aumentada $(A \mid B)$ para observar el comportamiento de los rangos:

\[
(A \mid B) = \left( \begin{array}{cccc|c} 
    1 & 2 & -1 & 1 & 1 \\ 
    2 & 4 & -2 & 2 & 2 \\ 
    1 & 3 & a & a-1 & 2 \\ 
    3 & 7 & a-2 & a+1 & 4 
\end{array} \right)
\xrightarrow{\substack{R_2 \to R_2 - 2R_1 \\ R_3 \to R_3 - R_1 \\ R_4 \to R_4 - 3R_1}}
\left( \begin{array}{cccc|c} 
    1 & 2 & -1 & 1 & 1 \\ 
    0 & 0 & 0 & 0 & 0 \\ 
    0 & 1 & a+1 & a-2 & 1 \\ 
    0 & 1 & a+1 & a-2 & 1 
\end{array} \right)
\]

\[
\xrightarrow{R_4 \to R_4 - R_3}
\left( \begin{array}{cccc|c} 
    1 & 2 & -1 & 1 & 1 \\ 
    0 & 1 & a+1 & a-2 & 1 \\ 
    0 & 0 & 0 & 0 & 0 \\ 
    0 & 0 & 0 & 0 & 0 
\end{array} \right)
\]

Como $\text{Rg}(A) = 2$ y $\text{Rg}(A|B) = 2$, el sistema es \textbf{compatible} para cualquier valor de $a$.
\begin{center}
    \fbox{$S = \{a \in \mathbb{R}\}$}
\end{center}

\textbf{b) Resolución del sistema para $a = 2$:}

Sustituimos $a = 2$ en la matriz escalonada obtenida:
\[ \left( \begin{array}{cccc|c} 
    1 & 2 & -1 & 1 & 1 \\ 
    0 & 1 & 3 & 0 & 1 
\end{array} \right) \]

De la segunda fila: $y + 3z = 1 \implies \mathbf{y = 1 - 3z}$. \\
Sustituimos en la primera fila: $x + 2(1 - 3z) - z + w = 1$:
\[ x + 2 - 6z - z + w = 1 \implies x - 7z + w = -1 \implies \mathbf{x = 7z - w - 1} \]

\begin{center}
\fbox{
\begin{minipage}{0.6\textwidth}
    \centering
    \textbf{Resultado final (Solución general):} \\
    \medskip
    El sistema tiene infinitas soluciones de la forma: \\
    $(x, y, z, w) = (7z - w - 1, 1 - 3z, z, w) \quad \forall z, w \in \mathbb{R}$
\end{minipage}
}
\end{center}
\end{solbox}
% --- EJERCICIO 13 ---
\begin{solbox}{Ejercicio 13}
    \begin{enumerate}
    \item \textbf{Análisis mediante el Teorema del Rango:}

Escalonamos la matriz aumentada $(A \mid B)$.
\[
(A \mid B) = \left( \begin{array}{ccc|c} 
    1 & 3 & \beta-1 & 4 \\ 
    \alpha & 9 & \alpha(\beta-1) & 4\alpha \\ 
    -1 & 6-3\alpha & 4-\beta & 2 \\ 
    2 & 6 & 2\beta-5 & 11-\alpha^2 
\end{array} \right)
\xrightarrow{\substack{R_2 \to R_2 - \alpha R_1 \\ R_3 \to R_3 + R_1 \\ R_4 \to R_4 - 2R_1}}
\left( \begin{array}{ccc|c} 
    1 & 3 & \beta-1 & 4 \\ 
    0 & 9-3\alpha & 0 & 0 \\ 
    0 & 9-3\alpha & 3 & 6 \\ 
    0 & 0 & -3 & 3-\alpha^2 
\end{array} \right)
\]
\[
\xrightarrow{R_3 \to R_3 - R_2}
\left( \begin{array}{ccc|c} 
    1 & 3 & \beta-1 & 4 \\ 
    0 & 9-3\alpha & 0 & 0 \\ 
    0 & 0 & 3 & 6 \\ 
    0 & 0 & -3 & 3-\alpha^2 
\end{array} \right)
\xrightarrow{R_4 \to R_4 + R_3}
\left( \begin{array}{ccc|c} 
    1 & 3 & \beta-1 & 4 \\ 
    0 & 9-3\alpha & 0 & 0 \\ 
    0 & 0 & 3 & 6 \\ 
    0 & 0 & 0 & 9-\alpha^2 
\end{array} \right)
\]

\item \textbf{Condiciones para infinitas soluciones:}
Para que el sistema tenga infinitas soluciones, requerimos que $\text{Rg}(A) = \text{Rg}(A|B) < 3$.
\begin{enumerate}
    \item Para que la fila 4 se anule: $9 - \alpha^2 = 0 \implies \alpha = 3$ o $\alpha = -3$.
    \item Si $\alpha = 3$, la fila 2 se convierte en $[0 \ 0 \ 0 \mid 0]$, lo que reduce el rango.
    \item Si $\alpha = -3$, la fila 2 queda $[0 \ 18 \ 0 \mid 0]$, manteniendo el rango en 3 (solución única).
\end{enumerate}

Si $\mathbf{\alpha = 3}$, la matriz queda con rango 2 independientemente de $\beta$, ya que la variable $\beta$ solo aparece en la primera fila multiplicando a $z$, lo cual no afecta la existencia de pivotes en este escalonamiento.
\end{enumerate}

\begin{center}
\fbox{
\begin{minipage}{0.4\textwidth}
    \centering
    \textbf{Resultado:} \\
    \medskip
    $S = \{(\alpha, \beta) \in \mathbb{R}^2 \mid \alpha = 3, \beta \in \mathbb{R}\}$
\end{minipage}
}
\end{center}
\end{solbox}
% --- EJERCICIO 14 ---
\begin{solbox}{Ejercicio 14}
    \textbf{1. Escalonamiento de la matriz aumentada $(A \mid B)$:}

\[
(A \mid B) =\left( \begin{array}{cccc|c} 
    1 & 2 & 3 & 0 & 2 \\ 
    3 & 6 & 1 & 8 & -3 \\ 
    5 & 10 & -1 & 17 & 12 
\end{array} \right)
\xrightarrow{\substack{R_2 \to R_2 - 3R_1 \\ R_3 \to R_3 - 5R_1}}
\left( \begin{array}{cccc|c} 
    1 & 2 & 3 & 0 & 2 \\ 
    0 & 0 & -8 & 8 & -9 \\ 
    0 & 0 & -16 & 17 & 2 
\end{array} \right)
\]

\[
\xrightarrow{R_3 \to R_3 - 2R_2}
\left( \begin{array}{cccc|c} 
    1 & 2 & 3 & 0 & 2 \\ 
    0 & 0 & -8 & 8 & -9 \\ 
    0 & 0 & 0 & 1 & 20 
\end{array} \right)
\]

\textbf{2. Análisis de Rangos e Incógnitas:}
\begin{itemize}
    \item El rango de la matriz de coeficientes es $\text{Rg}(A) = 3$.
    \item El rango de la matriz aumentada es $\text{Rg}(A \mid B) = 3$.
    \item Como $\text{Rg}(A) = \text{Rg}(A \mid B) < n$ (donde $n=4$), el sistema tiene \textbf{inrfinitas soluciones}.
\end{itemize}

\textbf{3. Determinación de variables:}
\begin{itemize}
    \item De la Fila 3: $\mathbf{x_4 = 20}$.
    \item De la Fila 2: $-8x_3 + 8(20) = -9 \implies -8x_3 + 160 = -9 \implies -8x_3 = -169 \implies \mathbf{x_3 = \frac{169}{8}}$.
    \item De la Fila 1: $x_1 + 2x_2 + 3(\frac{169}{8}) = 2 \implies x_1 + 2x_2 + \frac{507}{8} = \frac{16}{8} \implies \mathbf{x_1 = -2x_2 - \frac{491}{8}}$.
\end{itemize}

\begin{center}
\fbox{
\begin{minipage}{0.5\textwidth}
    \centering
    \textbf{Conjunto Solución:} \\
    \medskip
    $C.S. = \left\{ \left( -2\lambda - \frac{491}{8}, \lambda, \frac{169}{8}, 20 \right) \in \mathbb{R}^4 \mid \lambda \in \mathbb{R} \right\}$
\end{minipage}
}
\end{center}
\end{solbox}

\vspace{3cm}
\begin{center}
    \small\textbf{Usach Premium} — ¡Nos vemos en la ayudantía! \\
    \textit{"Por y para estudiantes."}
\end{center}

\end{soluciones}
\end{document}